\documentclass[11pt]{article}

\usepackage[margin=1in]{geometry}
\usepackage{amsmath,amssymb,amsthm}
\usepackage{array,booktabs}
\usepackage{enumitem}
\usepackage[hidelinks]{hyperref}

\theoremstyle{plain}
\newtheorem{proposition}{Proposition}
\newtheorem{lemma}{Lemma}
\newtheorem{corollary}{Corollary}
\newtheorem{conjecture}{Conjecture}
\theoremstyle{definition}
\newtheorem{remark}{Remark}

\newcommand{\F}{\mathbb{F}}
\newcommand{\Q}{\mathbb{Q}}
\newcommand{\N}{\mathbb{N}}
\newcommand{\ord}{\operatorname{ord}}
\newcommand{\Norm}{\operatorname{N}}
\newcommand{\Frob}{\operatorname{Frob}}
\newcommand{\nimmul}{\otimes}
\newcommand{\Qset}{\mathcal{Q}}
% claim-level tags (repo convention for this thread)
\newcommand{\PROVED}{\textbf{[\textsc{proved}]}}
\newcommand{\CERTIFIED}[1]{\textbf{[\textsc{certified} $#1$]}}
\newcommand{\CONSISTENT}{\textbf{[\textsc{consistent}]}}
\newcommand{\CONJECTURED}{\textbf{[\textsc{conjectured}]}}
\newcommand{\OPEN}{\textbf{[\textsc{open}]}}

\title{Draft: the finite excess in transfinite nim multiplication}
\author{a9lim}
\date{June 2026}

\begin{document}
\maketitle

\begin{abstract}
Conway's field $\mathrm{On}_2$ makes the ordinals below
$\omega^{\omega^{\omega}}$ an algebraic closure of $\F_2$ under
nim-arithmetic. Its multiplication is governed by a Kummer tower whose
carries are Lenstra's \emph{excess} values: for each odd prime $p$, the
carry is $\alpha_p=\kappa_{f(p)}+m_p$, where the transfinite part
$\kappa_{f(p)}$ has a known closed shape (Lenstra, DiMuro) but the finite
correction $m_p$ does not. This note consolidates the state of the excess
problem (\texttt{docs/OPEN.md} Problem~3 of \texttt{ogdoad}): the exact
finite-field reformulation of the root test as a multiplicative-order
criterion and its structural norm reduction; a complete analysis of the
$3$-power column via a half-angle splitting in the cyclotomic component
fields, certifying $m_r=1$ for all primes $r$ with
$f(r)\in\{3,9,\dots,3^6\}$ and forcing $m_p\geq4$ unconditionally on the
$f(p)=2\cdot3^k$ column --- now sharpened to $m_p=4$ \emph{exactly} at
every prime current factor tables reach (universally for $k\leq6$), via
a corrected compositum norm
$(\kappa+4)(\kappa+5)=\kappa^2+\kappa+\omega$ that collapses the $m=4$
test into the same trinomial field; the empirically exact $0/1/4$
candidate rule and the proof that the Lenstra component set alone cannot
decide the excess; and the reduction of Lenstra's open boundedness
question to that candidate.
Feasibility of the first open row, $p=719$, is assessed: a
Frobenius-orbit norm recurrence reduces it to arithmetic in
$\F_{2^{359}}$, with the cost concentrated in tower-aware Frobenius
composition. Each statement carries an explicit claim-level tag.
\end{abstract}

\paragraph{Claim levels.}
\PROVED{} -- unconditional, machine-checked or classically verified;
\CERTIFIED{k\leq6} -- proved for the stated range, blocked beyond it only
by unfactored cofactors; \CONSISTENT{} -- no counterexample found, not
proved; \CONJECTURED{} -- explicitly open in the literature or posed in
this thread; \OPEN{} -- neither confirmed nor refuted, no feasible attack
on record.

\section{Notation, implemented state, external data}\label{sec:notation}

\subsection{Notation}

For an odd prime $p$:
\begin{itemize}[leftmargin=*,itemsep=0pt]
\item $f(p)=\ord_p(2)$, the multiplicative order of $2$ modulo $p$;
\item $\kappa_h$ is the Lenstra/DiMuro tower element indexed by $h$;
\item $\Qset(h)$ is Lenstra's set of prime-power components appearing in
$\kappa_h$;
\item the \emph{excess} $m_p$ is the least finite $m$ such that
$\kappa_{f(p)}+m$ has no $p$-th root in the relevant finite component
field;
\item the Kummer carry is $\alpha_p=\kappa_{f(p)}+m_p$.
\end{itemize}
For the ordinal tower, a row with $\Qset(f(p))=\{q\}$ and finite excess
$m$ gives the ordinal sum corresponding to $\kappa_q+m$. Addition of
nimbers is XOR; $\nimmul$ is the nim product.

\subsection{Implemented tower state}

In the Rust tower (\path{src/scalar/big/ordinal/tower.rs}), as of 2026-06-13:
\begin{itemize}[leftmargin=*,itemsep=0pt]
\item the finite excess rows $m_p$ are now sourced from OEIS
A380496~\cite{OEIS}, the b-file's $126$ known rows (odd primes
$3 \le p \le 709$); the first $14$ rows reproduce DiMuro Table~1 + the
formerly locally verified $\alpha_{47}=\omega^{(\omega^7)}+1$;
\item the operational boundary is $\alpha_{719}$, the first OEIS-unknown
row: a carry needing $\alpha_{719}$ or beyond returns \texttt{None};
\item the table extends \emph{reach}, not \emph{feasibility}: each
$\alpha_p$ is reconstructed in code from $\ord_p(2)$, $\Qset(f(p))$, and
$m_p$, but for large $p$ the $\Qset$/finite-subfield reconstruction lives
in the degree-$e_p$ component field ($e_p$ in the millions; see the
snapshot below), so only the small-$e_p$ rows are materializable
end-to-end today.
\end{itemize}

\subsection{External data snapshot (2026-06-09)}

\begin{itemize}[leftmargin=*,itemsep=0pt]
\item OEIS A380496~\cite{OEIS} has $1417$ extended rows: $799$ known and
$618$ unknown; the b-file has $126$ initial known rows.
\item The first OEIS unknown is row $n=127$, the $127$th odd prime
$p=719$. For $p=719$: $f(719)=\ord_{719}(2)=359$ and
$\Qset(359)=\{359\}$.
\item The transfinite-nim-calculator logs~\cite{Peeters} record the direct
component exponent as $e_{719}=1{,}258{,}230{,}380$, the practical wall
for direct exponentiation.
\end{itemize}

\section{The order criterion and structural norm reduction}
\label{sec:norm}

\subsection{The exact criterion}

Throughout, $E$ denotes the degree of the component field containing
$\beta=\kappa_{f(p)}+m$, so $f=f(p)\mid E$ and hence
$p\mid 2^{E}-1$.

\begin{proposition}[power criterion]\label{prop:power}
\PROVED{} $\beta\in\F_{2^E}^{\times}$ has no $p$-th root in $\F_{2^E}$
iff $\beta^{(2^E-1)/p}\neq1$. Hence
\[
  m_p=\min\bigl\{m\geq0:\ (\kappa_{f(p)}+m)^{(2^E-1)/p}\neq1\bigr\}.
\]
\end{proposition}

\begin{proof}
$\F_{2^E}^{\times}$ is cyclic of order $2^E-1$ and $p\mid 2^E-1$; in a
cyclic group of order $N$, $\beta$ is a $p$-th power iff
$\beta^{N/p}=1$.
\end{proof}

\begin{corollary}[order form]\label{cor:order}
\PROVED{} If $v_p(2^E-1)=1$, then $\beta$ has no $p$-th root in
$\F_{2^E}$ iff $p\mid\ord(\beta)$, so
$m_p$ is the least $m$ with $p\mid\ord(\kappa_{f(p)}+m)$.
\end{corollary}

\begin{remark}[Wieferich caveat]\label{rem:wieferich}
In general ``no $p$-th root'' is equivalent to
$v_p(\ord\beta)=v_p(2^E-1)$, which collapses to
$p\mid\ord(\beta)$ only when $v_p(2^E-1)=1$. Since
$v_p(2^E-1)=v_p(2^{f(p)}-1)$ whenever $p\nmid E/f(p)$ (lifting the
exponent), the hypothesis can fail only at base-$2$ Wieferich-type primes
with $p^2\mid 2^{f(p)}-1$; the only known instances are $p=1093$ and
$p=3511$, both of which lie inside the extended A380496 range. For those
rows the exact test of Proposition~\ref{prop:power} (or the full-$p$-part
order condition) must be used. The certifications of
Section~\ref{sec:family} compute orders exactly, with squarefreeness
checks, and are unaffected. The research-note predecessor of this
document stated only the order form; the power criterion is the
correct general statement.
\end{remark}

\subsection{Norm reduction}

\begin{proposition}[norm reduction]\label{prop:norm}
\PROVED{} For $\beta\in\F_{2^E}^{\times}$ with $f=f(p)\mid E$,
\[
  \beta^{(2^E-1)/p}
  =
  \Norm_{\F_{2^E}/\F_{2^f}}(\beta)^{(2^f-1)/p}.
\]
\end{proposition}

\begin{proof}
$\Norm_{\F_{2^E}/\F_{2^f}}(\beta)=\beta^{(2^E-1)/(2^f-1)}$, and
$\bigl((2^E-1)/(2^f-1)\bigr)\cdot\bigl((2^f-1)/p\bigr)=(2^E-1)/p$.
\end{proof}

So the $p=719$ test reduces to: (1) compute
$N=\Norm_{\F_{2^E}/\F_{2^{359}}}(\kappa_{359}+1)$ in $\F_{2^{359}}$;
(2) test whether $719\mid\ord(N)$ in $\F_{2^{359}}^{\times}$, a field of
degree $359$ in which $719\mid2^{359}-1$.

\begin{proposition}[Galois-conjugacy invariance]\label{prop:galois}
\PROVED{} The order of $\kappa_{f(p)}+m$ is constant on its Frobenius
orbit in $\F_{2^E}$, and the norm
$N=\prod_{i=0}^{E/f-1}\Frob^{if}(\kappa_{f}+m)$ is determined
structurally by the tower equations without materialising the full orbit.
\end{proposition}

\begin{proposition}[tower pinning]\label{prop:pinning}
\PROVED{} The tower equations pin $\kappa_{359}$ up to Frobenius
conjugacy: each composed minimal polynomial along the component chain
$359\to179\to89\to11\to5\to\text{finite}$ is irreducible (the
non-$q$-th-power property defining the lower excess guarantees this), so
$\kappa_{359}$ is a well-defined conjugacy class and $N$ is a well-defined
element of $\F_{2^{359}}$.
\end{proposition}

\subsection{Feasibility for $p=719$}

The component field degree is
$E=2\cdot2\cdot5\cdot11\cdot89\cdot179\cdot359=1{,}258{,}230{,}380$; the
norm requires $E/f=3{,}504{,}820$ Frobenius steps.

\CONSISTENT{} The norm can in principle be computed by the
binary-splitting Frobenius recurrence
$N_{2k}=N_k\cdot\Frob^{fk}(N_k)$ (iterated-Frobenius technique, von zur
Gathen--Shoup~\cite{vzGS92}, Kaltofen--Shoup~\cite{KS98}), requiring about
$\log_2(3{,}504{,}820)\approx22$ squarings-with-Frobenius-composition at
degree $E$ over $\F_2$. Each Frobenius composition is a modular
composition in $\F_2[x]$ modulo a degree-$E$ polynomial; with dense
polynomial arithmetic at $E\approx1.26\times10^9$ this is not locally
feasible.

\CONJECTURED{} (as a cost model, not a theorem) Kummer-tower-aware
Frobenius -- standard-lattice arithmetic for compatibly embedded finite
fields in the style of De~Feo--Randriam--Rousseau~\cite{DFRR19}, where
$q$-power Frobenius acts near-monomially per tower level -- is the likely
$10$--$100\times$ lever: the sparse factored structure
$E=2\cdot2\cdot5\cdot11\cdot89\cdot179\cdot359$ admits a stepwise
Frobenius representation in which each level contributes a small-degree
modular composition. A \emph{measured} cost model, not a run, is the
correct first deliverable.

\OPEN{} Whether the Kummer-tower Frobenius model makes the $m_{719}$
computation feasible in wall-clock hours rather than months on available
hardware (a single consumer GPU; GF(2) arithmetic is XOR-sliced and
GPU-friendly).

\begin{remark}[shared abstraction]\label{rem:bridgek}
The structural primitive needed here --
$\Norm_{E/K}(\beta)=\prod_i\Frob^{i}(\beta)$ computed by a
Frobenius-orbit recurrence -- is the same primitive that the cyclic-algebra
Brauer-invariant bridge (Bridge~K, \path{docs/TBD.md}) requires for
reduced norms over a Frobenius-generated cyclic extension. The existing
\texttt{FieldExtension::relative\_norm} does not apply to the
term-algebra/transfinite setting, but the shape is identical; factoring
out a reusable \path{relative_norm_over_frobenius_orbit} is an
engineering opportunity. This is not a claim that the bounded
\texttt{Fpn} norm certifies $m_{719}$.
\end{remark}

\section{The $3^k$ family}\label{sec:family}

Write $\zeta:=\kappa_{3^k}$ and $h:=3^k$ throughout this section.

\subsection{Structure}

\begin{proposition}[cyclotomic recognition]\label{prop:cyclo}
\PROVED{} $\zeta^{3^k}=\kappa_2$, which has order $3$; hence $\zeta$ is a
primitive $3^{k+1}$-st root of unity. Since $2$ is a primitive root
modulo $3^{k+1}$, the cyclotomic polynomial
$\Phi_{3^{k+1}}(x)=x^{2h}+x^h+1$ is irreducible over $\F_2$, the
component field is $\F_2(\zeta)=\F_{2^{2h}}$, and it carries the index-$2$
subfield $L=\F_{2^h}$.
\end{proposition}

\begin{proposition}[half-angle splitting]\label{prop:halfangle}
\PROVED{} With $s=(3^{k+1}+1)/2$ (the inverse of $2$ modulo $3^{k+1}$),
\[
  \kappa_{3^k}+1=\zeta^{s}\,\bigl(\zeta^{s}+\zeta^{-s}\bigr),
\]
where $\zeta^{s}$ lies in the norm-one circle $U$ (of order exactly
$3^{k+1}$ within the subgroup of order $2^h+1$) and
$(\zeta^{s}+\zeta^{-s})^2=\zeta+\zeta^{-1}=:\gamma_k\in L^{\times}$.
Since $\F_{2^{2h}}^{\times}=L^{\times}\times U$ with coprime orders,
\[
  \ord(\kappa_{3^k}+1)=3^{k+1}\cdot\ord(\gamma_k),
  \qquad
  \ord(\gamma_k)\mid 2^{3^k}-1.
\]
Machine-checked corollaries: $(\kappa+1)^{2^h-1}=\zeta^{-1}$ and
$\Norm_{F/L}(\kappa+1)=\gamma_k$.
\end{proposition}

\begin{proposition}[translates $m=2,3$]\label{prop:translates}
\PROVED{} $\kappa_{3^k}+2$ and $\kappa_{3^k}+3$ split as in
Proposition~\ref{prop:halfangle} with $\gamma_k$ replaced by a Galois
conjugate. Therefore
$\ord(\kappa_{3^k}+m)=3^{k+1}\cdot\ord(\gamma_k)$ for $m=1,2,3$, and
$\ord(\kappa_{3^k})=3^{k+1}$.
\end{proposition}

\begin{proposition}[the $2\cdot3^k$ exception, unconditionally]
\label{prop:exception}
\PROVED{} Any prime $p$ with $f(p)=2\cdot3^k$ satisfies
$p\mid 2^{3^k}+1$, hence $p$ divides neither $3^{k+1}$ nor $2^{3^k}-1$,
hence $p$ never divides $\ord(\kappa_{3^k}+m)$ for $m\in\{0,1,2,3\}$. By
the order criterion, $m_p\geq4$. The full primitivity conjecture $C_k$
below is not needed; the splitting alone forces the bound. The argument
applies beyond the current tables, e.g.\ $p=87211$, $f=54$.
\end{proposition}

\begin{proposition}[the $3$-power column]\label{prop:column}
\PROVED{} For a prime $r$ with $f(r)=3^k$, the values $m\in\{0,2,3\}$ are
impossible: $m_r=0$ would require $r\mid\ord(\kappa_{3^k})=3^{k+1}$, a
power of $3$, while $r\equiv1\pmod{3^k}$ is a prime distinct from $3$;
and $m=2,3$ share $m=1$'s order class by
Proposition~\ref{prop:translates}. Any failure of $C_k$ therefore gives
$m_r\geq4$ directly.
\end{proposition}

\begin{proposition}[norm tower]\label{prop:tower}
\PROVED{} $\gamma_{k-1}=\gamma_k^3+\gamma_k
=\Norm_{L_k/L_{k-1}}(\gamma_k)$ (minimal polynomial
$X^3+X+\gamma_{k-1}$), so $\ord(\gamma_{k-1})\mid\ord(\gamma_k)$. With
lifting-the-exponent ($v_r(2^{3^k}-1)=v_r(2^{f(r)}-1)$ for old primes
$r\neq3$), full-order parts propagate upward, and conjecture $C_k$
reduces to $C_{k-1}$ plus primitivity at the new primes
$r\mid\Phi_{3^k}(2)$.
\end{proposition}

\begin{conjecture}[$C_k$: $\gamma$-primitivity]\label{conj:ck}
$\gamma_k$ is primitive in $L^{\times}=\F_{2^{3^k}}^{\times}$.
\end{conjecture}

\begin{proposition}[equivalence]\label{prop:equiv}
\PROVED{} For primes $r$ with $f(r)\mid 3^k$ a $3$-power:
\[
  C_k
  \iff
  m_r=1\ \text{for every prime $r$ with } f(r)\in\{3,9,\dots,3^k\}.
\]
The family formula is the $0/1/4$ candidate of Section~\ref{sec:rule}
restricted to the $3$-power column.
\end{proposition}

\subsection{Certification status}

\CERTIFIED{k\leq6} $C_k$ holds for $k=1,\dots,6$: $\gamma_k$ is primitive
in $\F_{2^{3^k}}^{\times}$, and
$\ord(\kappa_{3^k}+1)=3^{k+1}\,(2^{3^k}-1)$ is verified exactly --
product reconstruction, squarefreeness, primality by deterministic
Miller--Rabin below $3.3\times10^{24}$ and MR-64 probable-primality above,
an independent sieve re-deriving the small factors, cross-validation
against the independent term algebra for $k=1,2$ and against the
calculator rows $m_{2593}=1$, $m_{487}=1$
(\path{experiments/cyclotomic_3k_family.py},
\path{experiments/excess/}).

\CONSISTENT{} $k=7,8$: every known prime factor of $\Phi_{3^7}(2)$ and
$\Phi_{3^8}(2)$ divides $\ord(\gamma_k)$. Full certification is blocked
only by the unfactored cofactors (FactorDB status composite-no-factor);
no failure has been found.

\OPEN{} Whether ECM or GNFS can factor the $\Phi_{3^7}(2)$ and
$\Phi_{3^8}(2)$ cofactors within a realistic budget, converting $k=7,8$
from \textsc{consistent} to \textsc{certified}.

\subsection{Certified excess rows produced}

The following $m_r=1$ rows are certified by the order criterion
independently of the cofactor issue (analysis-level results; they are not
new Rust carries):
\begin{center}
\begin{tabular}{rl}
\toprule
$f$ & primes $r$ with $m_r=1$ \\
\midrule
$27$   & $262657$ \\
$81$   & $71119,\ 97685839$ \\
$243$  & $16753783618801,\ 192971705688577,\ 3712990163251158343$ \\
$729$  & $80191,\ 97687,\ 379081,\ \mathrm{P}42,\ \mathrm{P}90$ \\
$2187$ & $39367,\ 7606246033,\ 263196614521,\ 529063556041$ \\
$6561$ & $209953,\ 1299079,\ 70063267397606709277393$ \\
\bottomrule
\end{tabular}
\end{center}
($\mathrm{P}42$, $\mathrm{P}90$ denote the certified/probable primes of
$42$ and $90$ decimal digits dividing $\Phi_{3^6}(2)$.)

\section{The $0/1/4$ candidate}\label{sec:rule}

\begin{conjecture}[the $0/1/4$ rule]\label{conj:rule}
\CONSISTENT{} (not proved)
\[
  m_p=
  \begin{cases}
    0 & \text{if $\Qset(f(p))$ is not a singleton odd prime-power,}\\
    4 & \text{if $f(p)=2\cdot3^k$, $k\geq1$,}\\
    1 & \text{otherwise (singleton odd prime-power $\Qset(f(p))$).}
  \end{cases}
\]
The rule matches all $950$ calculator records with known $\Qset$-sets and
every OEIS-known row covered by those $\Qset$-sets.
\end{conjecture}

\subsection{The component set alone is insufficient}

\PROVED{} $\Qset(f(p))$ alone does not determine $m_p$:
\begin{center}
\begin{tabular}{ccc}
\toprule
$\Qset$ & $m=4$ & $m=1$ \\
\midrule
$\{9\}$   & $p=19$,   $f=18$  & $p=73$,   $f=6$ \\
$\{81\}$  & $p=163$,  $f=162$ & $p=2593$, $f=18$ \\
$\{243\}$ & $p=1459$, $f=486$ & $p=487$,  $f=18$ \\
\bottomrule
\end{tabular}
\end{center}

\PROVED{} The split for $\Qset=\{9\}$ is explained by the order
criterion: $\ord(\kappa_9+1)=3^3\cdot(2^9-1)=27\cdot511=13797$ and
$511=7\cdot73$, so $73\mid\ord(\kappa_9+1)$ but
$19\nmid\ord(\kappa_9+1)$; adding $m=4$ changes the order and introduces
the factor $19$.

\PROVED{} The lower bound of the exception column is unconditional:
$m_p\geq4$ for all $p$ with $f(p)=2\cdot3^k$
(Proposition~\ref{prop:exception}).

\CONJECTURED{} The matching upper bound $m_p\leq4$ on that column;
equivalent to $C_k$ for the primes in the exception column.

\subsection{The $m=4$ upper bound: settled at every visible prime}
\label{sec:m4}

The $m=4$ translate $\kappa_{3^k}+4$ involves the nimber $4\in\F_{16}$,
which does not lie in $\F_2(\zeta)$ (as $4\nmid2\cdot3^k$); the
translate lives in the degree-$4\cdot3^k$ compositum $\F_{2^{4h}}$,
$h=3^k$. (June 2026 pass; probe
\path{experiments/exception_column_m4.py}.)

\begin{proposition}[the corrected norm]\label{prop:m4norm}
\PROVED{} Let $\sigma=\Frob^{2h}$ generate
$\mathrm{Gal}(\F_{2^{4h}}/\F_{2^{2h}})$. Then $\sigma(4)=5$, and with
$\omega:=\zeta^{h}=\kappa_2$,
\[
  N:=\Norm_{\F_{2^{4h}}/\F_{2^{2h}}}(\kappa_{3^k}+4)
   =(\kappa+4)(\kappa+5)
   =\kappa^2+\kappa+\omega .
\]
\end{proposition}

\begin{proof}
$\F_{16}\cap\F_{2^{2h}}=\F_4$ because $\gcd(4,2\cdot3^k)=2$, so
$\sigma|_{\F_{16}}$ is the nontrivial element of
$\mathrm{Gal}(\F_{16}/\F_4)$, which swaps the two roots $4,5$ of the
Artin--Schreier minimal polynomial $y^2+y+\omega$ of $4$ over $\F_4$.
(Equivalently: $2h\equiv2\pmod4$ and the Frobenius orbit of $4$ is
$4\to6\to5\to7$, so $\sigma(4)=4^{2^2}=5$.) Then
$(\kappa+4)(\kappa+5)=\kappa^2+\kappa+(4\nimmul5)$ and
$4\nimmul5=4^2+4=\omega$ by the same Artin--Schreier relation.
\end{proof}

\begin{remark}[erratum]\label{rem:erratum}
The 2026-06-10 draft of this note stated the norm as
$(\kappa+4)(\kappa+\Frob^2(4))=(\kappa+4)(\kappa+6)$. The conjugate
$6=4^2$ is $\Frob^1(4)$, one squaring short of $\Frob^2(4)=5$. The slip
is visible without computation: $4\nimmul6=14\notin\F_4$, so
$(\kappa+4)(\kappa+6)\notin\F_{2^{2h}}$ and cannot be a norm.
\end{remark}

The corrected norm collapses the $m=4$ test into the \emph{same} sparse
trinomial model $\F_2[x]/(x^{2h}+x^h+1)$ as the rest of the $3$-power
analysis: no compositum arithmetic is needed.

\begin{corollary}[in-field criterion]\label{cor:m4crit}
\PROVED{} For a prime $p$ with $f(p)=2\cdot3^k$ and $v_p(2^{2h}-1)=1$,
write $\overline{N}=N^{2^h}$ and $M_k:=\overline{N}/N=N^{2^h-1}\in U$.
Then
\[
  m_p=4
  \iff
  p\mid\ord(M_k)
  \iff
  M_k^{(2^h+1)/p}\neq1 .
\]
\end{corollary}

\begin{proof}
$m_p\geq4$ is Proposition~\ref{prop:exception}, so $m_p=4$ iff
$\kappa+4$ has no $p$-th root. By Propositions~\ref{prop:power}
and~\ref{prop:norm} with $E=4h$, $f=2h$ (and
$v_p(2^{4h}-1)=v_p(2^{2h}-1)=1$ by lifting the exponent, since
$p\nmid2$), that holds iff $N^{(2^{2h}-1)/p}\neq1$. Since
$\F_{2^{2h}}^{\times}=L^{\times}\times U$ with coprime orders,
$p\mid2^h+1$, and $(2^{2h}-1)/p=(2^h-1)\cdot\bigl((2^h+1)/p\bigr)$, the
test equals $M_k^{(2^h+1)/p}\neq1$.
\end{proof}

\begin{conjecture}[$D_k$: the column analogue of $C_k$]\label{conj:dk}
The prime-to-$3$ part of $\ord(M_k)$ is full:
$\bigl((2^h+1)/3^{k+1}\bigr)\mid\ord(M_k)$. Equivalently (given the
per-prime squarefreeness audits): $m_p=4$ for \emph{every} prime $p$
with $f(p)=2\cdot3^k$.
\end{conjecture}

\CERTIFIED{k\leq6} $D_k$ holds for $k=2,\dots,6$: every prime factor of
the fully factored $\Phi_{2\cdot3^k}(2)$ divides $\ord(M_k)$. Hence
$m_p=4$ exactly --- not merely $\geq4$ --- for every prime $p$ with
$f(p)\in\{18,54,162,486,1458\}$:
\begin{center}
\begin{tabular}{rl}
\toprule
$f$ & primes $p$ with $m_p=4$ \\
\midrule
$18$   & $19$ \\
$54$   & $87211$ \\
$162$  & $163,\ 135433,\ 272010961$ \\
$486$  & $1459,\ 139483,\ 10429407431911334611,$ \\
       & $918125051602568899753$ \\
$1458$ & $227862073,\ 3110690934667,\ 216892513252489863991753,$ \\
       & $1102099161075964924744009,\ \mathrm{P}78$ \\
\bottomrule
\end{tabular}
\end{center}
($\mathrm{P}78$ is the certified $78$-digit prime of
$\Phi_{2\cdot3^6}(2)$; PRP-local here, factordb marks it proven.) The
anchor rows $m_{19}=m_{163}=m_{1459}=4$ (DiMuro / calculator) are
reproduced, never assumed; the remaining eleven rows are new.
Independent cross-checks: the term algebra of
\path{experiments/ordinal_excess_probe.py} re-derives the $p=19$ root
pattern ($m=0,\dots,3$ root, $m=4$ no root) and $m_{87211}=4$; an
explicit Artin--Schreier compositum model $\F[y]/(y^2+y+\omega)$
verifies $\sigma(4)=5$, the norm identity, and the full direct power
test $(\kappa+4)^{(2^{4h}-1)/p}\neq1$ on the smallest prime of each
level $k\leq6$.

\CONSISTENT{} $k=7,8$: every \emph{known} prime factor of
$\Phi_{2\cdot3^7}(2)$ and $\Phi_{2\cdot3^8}(2)$ --- seven and five
primes respectively; factordb CF entries verified locally by
divisibility, primality, and order tests, the small ones re-derived by
sieve, and the list cross-audited against the raw factordb API after an
adversarial review caught a dropped $43$-digit $k=7$ prime --- also has
$m_p=4$. An $m_p\geq5$ example, if one exists, now hides strictly
inside the unfactored cofactors.

\begin{proposition}[twisted norm tower]\label{prop:twisted}
\PROVED{} With $\eta=\zeta^3=\kappa_{3^{k-1}}$,
\[
  \Norm_{\F_{2^{2\cdot3^k}}/\F_{2^{2\cdot3^{k-1}}}}(N_k)
  =\eta^2+\omega^2\eta+1 ,
\]
which is neither $N_{k-1}=\eta^2+\eta+\omega$, nor one of its Frobenius
conjugates (their exponent pattern is $(2^{j+1},2^j)$), nor a scaling
$N_{k-1}(\lambda\eta)$ with $\lambda\in\mu_3$.
\end{proposition}

\begin{proof}
The generator $\tau$ of the degree-$3$ step sends $\zeta\mapsto
\zeta^{c}$ with $c=2^{2\cdot3^{k-1}}$ a primitive cube root of unity
modulo $3^{k+1}$; the cube roots are $1+a3^k$, so
$\tau(\zeta)=\omega^{a}\zeta$ ($a\neq0$) while $\tau(\omega)=\omega$.
Hence the norm is $\prod_{r^3=\eta}(r^2+r+\omega)$ over the three cube
roots $r$ of $\eta$. Writing $\alpha,\beta$ for the two roots of
$T^2+T+\omega$ (so $\alpha+\beta=1$, $\alpha\beta=\omega$), each factor
is $(r+\alpha)(r+\beta)$, and the product collapses by the resolvent
\[
  \prod_{r^3=\eta}(r^2+r+\omega)
  =(\alpha^3+\eta)(\beta^3+\eta)
  =\eta^2+(\alpha^3+\beta^3)\eta+(\alpha\beta)^3
  =\eta^2+\omega^2\eta+1,
\]
since $\alpha^3+\beta^3=(\alpha+\beta)^3+\alpha\beta(\alpha+\beta)
=1+\omega=\omega^2$ and $(\alpha\beta)^3=\omega^3=1$. Machine-checked at
$k=2,3,4$.
\end{proof}

So, unlike the $\gamma$ tower of Proposition~\ref{prop:tower},
old-prime full-order parts do \emph{not} propagate up the column, and
$D_k$ is genuinely per-level: the certification at each $k$ stands on
its own.

\begin{remark}[the $3$-part]\label{rem:threepart}
\CONSISTENT{} The $3$-part of $\ord(M_k)$ carries no column information
but is strikingly regular in range: $v_3(\ord M_k)=k+1$ (full) for odd
$k$ and $k-1$ (deficiency exactly $3^2$) for even $k$, for
$2\leq k\leq6$. No explanation is on record.
\end{remark}

\section{Boundedness}\label{sec:bounded}

\CONJECTURED{} (Lenstra~\cite{Lenstra77}) $m_p$ is globally bounded.
Lenstra gave unconditional lower-bound rules (singleton odd $\Qset$
forces positive excess; $f(p)=2\cdot3^k$ forces excess at least $4$) but
left absolute boundedness open.

\PROVED{} If the $0/1/4$ candidate holds, then $m_p\leq4$ for all $p$;
proving Conjecture~\ref{conj:rule} would settle boundedness.

\PROVED{} The $3$-power column cannot produce $m\in\{2,3\}$
(Proposition~\ref{prop:column}); any $C_k$ failure jumps directly to
$m_r\geq4$.

\PROVED{} (conditional on $C_k$) For every prime $r$ with $f(r)=3^j$,
$j\leq k$: $C_k$ implies $m_r=1$. With $C_k$ certified for $k\leq6$, all
primes in the $3$-power column with $f(r)\leq6561$ have $m_r=1$
unconditionally.

\CERTIFIED{k\leq6} (June 2026 pass) On the $2\cdot3^k$ exception column
the matching upper bound is now certified at every visible prime:
$m_p=4$ exactly for every prime with $f(p)=2\cdot3^k$, $k\leq6$, and for
every known prime at $k=7,8$ (Section~\ref{sec:m4}). Boundedness on the
exception column therefore holds unconditionally through $f=1458$ and is
equivalent to $D_k$ (Conjecture~\ref{conj:dk}) beyond.

\OPEN{} Boundedness outside the $3$-power and $2\cdot3^k$ columns. The
non-cyclotomic singleton chains (e.g.\ the $11$-chain and the $23$, $29$,
$47$ components) have $m=1$ rows in the calculator data and OEIS, but no
order-criterion proof covers them the way $C_k$ covers the $3$-power
column.

\OPEN{} Whether any prime $p$ has $m_p\geq5$. No counterexample is known;
the question is not known to be decidable by any current structural
argument outside the $3$-adic columns.

\section{Ranked next moves}\label{sec:moves}

\begin{enumerate}[leftmargin=*]
\item \emph{Dress rehearsal on $m_{89}$, $m_{179}$, $m_{359}$ (highest
priority).} Certify by direct order tests the calculator rows on which a
$m_{719}$ certificate depends. The component fields have degrees
$E=19{,}580$ ($p=89$) and $E=3{,}504{,}820$ ($p=359$ -- a stretch: hours
to a day with \texttt{gf2x}-style arithmetic). This calibrates the cost
model for the next move on the same code path and settles the provenance
question (``what evidence is acceptable'') before any new $\alpha$ ships.
Confirm $m_{359}=1$ first; the norm recurrence for $p=719$ shifts if
not. \emph{Feasible now; blocks the next move.}
\item \emph{Certify $m_{719}$ via the Kummer-tower Frobenius norm.}
Implement the binary-splitting norm $N_{2k}=N_k\cdot\Frob^{359k}(N_k)$ in
the tower-sparse representation; measure the cost model on the calibrated
path; if tower-aware Frobenius brings the ${\sim}44$ modular compositions
to feasible size, run to certify or falsify $m_{719}=1$. Factor out the
shared \texttt{relative\_norm\_over\_frobenius\_orbit} abstraction
(Remark~\ref{rem:bridgek}). \emph{Blocked on move 1 for the cost
model.}
\item \emph{[Done, 2026-06-12.] $m_p=4$ exactly on the $2\cdot3^k$
column.} Settled at every prime current factor tables reach
(Section~\ref{sec:m4}, \path{experiments/exception_column_m4.py}): the
corrected norm $(\kappa+4)(\kappa+5)=\kappa^2+\kappa+\omega$ collapses
the test in-field; certified universally for $k\leq6$, consistent at
$k=7,8$. Successor moves: prove $D_k$ (Conjecture~\ref{conj:dk}) ---
the twisted norm tower (Proposition~\ref{prop:twisted}) shows why
$\gamma$-style propagation is unavailable; factor the
$\Phi_{2\cdot3^7}(2)$ and $\Phi_{2\cdot3^8}(2)$ cofactors (joins the
ECM/GNFS batch of the next move); explain the $3$-part parity
(Remark~\ref{rem:threepart}).
\item \emph{Convert $C_7$, $C_8$ from \textsc{consistent} to
\textsc{certified}.} ECM/GNFS jobs on the cofactors of $\Phi_{3^7}(2)$
and $\Phi_{3^8}(2)$; full factorization also emits a batch of new
A380496 rows. \emph{Background job; uncertain outcome; blocks
nothing.}
\item \emph{Exact-order census on the non-cyclotomic singleton chains.}
Compute $\ord(\kappa_q+m)$ for small $m$ in the explicit component fields
for the $11$-chain and the $23$, $29$, $47$ components (degrees up to
${\sim}5060$; trivial locally); look for an analogue of the
$L^{\times}\times U$ splitting and for any forced-translate mechanism.
This is the main structural gap for boundedness outside the $3$-power
columns. \emph{Feasible now; no structural theory yet.}
\end{enumerate}

\section{Provenance and validation}\label{sec:provenance}

\subsection{Independent oracle}

\path{experiments/ordinal_excess_probe.py} is a small local term-algebra
oracle -- not a replacement for CGSuite~\cite{CGSuite} or the C++
calculator~\cite{Peeters}, but a way to verify the first subtle cases
without using the Rust production tower as its own oracle. It implements
the impartial term algebra used by the calculators, a
multiplicative-order test for small component fields, and a fixed-base
exponentiation path (ported from the calculator's strategy) for targeted
root tests where full order factorization is unnecessary. Current probe
output:

\begin{center}
\begin{tabular}{rrlc}
\toprule
$p$ & $m$ & $\Qset$ & root? \\
\midrule
$7$  & $0$ & $\{3\}$  & yes \\
$7$  & $1$ & $\{3\}$  & no \\
$19$ & $1$ & $\{9\}$  & yes \\
$19$ & $4$ & $\{9\}$  & no \\
$73$ & $1$ & $\{9\}$  & no \\
$47$ & $1$ & $\{23\}$ & no \\
\bottomrule
\end{tabular}
\end{center}

The last line certifies $m_{47}=1$ using only lower verified rows; since
$f(47)=23$ and $\Qset(23)=\{23\}$, the carry is
$\alpha_{47}=\kappa_{23}+1=\omega^{(\omega^7)}+1$, the value now
implemented in \path{tower.rs}.

\subsection{The $\alpha_{47}$ promotion}

\emph{Superseded 2026-06-13:} this pass moved the boundary to
$\alpha_{53}$ by promoting a single locally verified row; the finite
excess table is now the full OEIS A380496 b-file ($126$ rows, odd primes
$3 \le p \le 709$), so the refusal boundary sits at $\alpha_{719}$ and the
boundary test is \path{boundary_returns_none_past_prime_709}. The original
pass is recorded below for provenance.

The shipped state at the time: \path{src/scalar/big/ordinal/tower.rs} carries
$\alpha_{47}$ with the landmark test
\path{locally_verified_alpha_47_landmark} and moves the refusal
boundary to $\alpha_{53}$; the documented boundaries in
\path{nim.rs}, \path{mod.rs}, \path{src/games/nimber_game.rs},
\texttt{README.md}, and \texttt{docs/OPEN.md} were updated in step. The full
gate ran clean: \texttt{cargo fmt --check}, \texttt{cargo test},
\texttt{cargo check}/\texttt{clippy} with and without the
\texttt{python} feature (warnings denied), and the probe under
\texttt{py\_compile}.

\subsection{Supporting probes}

The $3^k$-family certification chain is
\path{experiments/cyclotomic_3k_family.py} (committed, maintained),
joined 2026-06-12 by \path{experiments/exception_column_m4.py} (the
$2\cdot3^k$ exception column, Section~\ref{sec:m4}; committed,
maintained, stdlib-only, $\sim$2 minutes), with the research-run probes
under \path{experiments/excess/} (machine-generated, rescued from
ephemeral storage; triage before citing). The analysis in
Sections~\ref{sec:norm}--\ref{sec:bounded} synthesizes the 2026-06-10
parallel research run plus the 2026-06-12 exception-column pass; each
claim above carries its own tag, and no claim depends on an unverified
solver.

\section*{Acknowledgements}
The \texttt{ogdoad} library, the experiments, and this draft were
developed with LLM assistance; the analysis consolidates a structured
parallel research run (2026-06-10). Mathematical direction, framing, and
responsibility for claims remain the author's.

\begin{thebibliography}{9}
\bibitem{Conway76} J.~H.~Conway, \emph{On Numbers and Games}, Academic
Press, 1976.
\bibitem{DFRR19} L.~De~Feo, H.~Randriam, \'E.~Rousseau, \emph{Standard
lattices of compatibly embedded finite fields}, Proc.\ ISSAC 2019;
arXiv:1906.00870.
\bibitem{DiMuro} J.~DiMuro, \emph{On $\mathrm{On}_p$}, arXiv:1108.0962.
\bibitem{KS98} E.~Kaltofen, V.~Shoup, \emph{Subquadratic-time factoring of
polynomials over finite fields}, Math.\ Comp.\ \textbf{67} (1998),
1179--1197.
\bibitem{Lenstra77} H.~W.~Lenstra, Jr., \emph{On the algebraic closure of
two}, Indag.\ Math.\ \textbf{39} (1977), 389--396;
\url{https://pub.math.leidenuniv.nl/~lenstrahw/PUBLICATIONS/1977e/art.pdf}.
\bibitem{Lenstra78} H.~W.~Lenstra, Jr., \emph{Nim multiplication},
S\'eminaire de Th\'eorie des Nombres de Bordeaux, 1977--1978.
\bibitem{OEIS} OEIS Foundation Inc., entry A380496,
\emph{The On-Line Encyclopedia of Integer Sequences},
\url{https://oeis.org/A380496}.
\bibitem{CGSuite} A.~N.~Siegel, \emph{CGSuite} (software),
\texttt{NimFieldCalculator},
\url{https://github.com/aaron-siegel/cgsuite}.
\bibitem{Peeters} D.~Peeters, \emph{transfinite-nim-calculator}
(software),
\url{https://github.com/DjangoPeeters/transfinite-nim-calculator}.
\bibitem{vzGS92} J.~von~zur~Gathen, V.~Shoup, \emph{Computing Frobenius
maps and factoring polynomials}, Comput.\ Complexity \textbf{2} (1992),
187--224.
\end{thebibliography}

\end{document}
